\subsection{Expert Générateur de peptides}
\begin{frame}
    \frametitle{Générateur}
    \begin{block}{Modèle du générateur}
        \begin{itemize}
        \item Modèle générant des séquences de peptides à partir de séquence ou propriétés souhaitées
        \item Entraîné sur une base de données de peptides connus et leurs propriétés
        \item Capable de générer des peptides avec des propriétés spécifiques (ex: propriétés chimiques)
        \end{itemize}
    \end{block}

    \begin{block}{Modèles implémentés}
        \begin{itemize}
        \item VAE (Variational Autoencoder): apprend une représentation latente des peptides et génère de nouvelles séquences à partir de cette représentation
        \item CVAE (Conditional Variational Autoencoder): extension du VAE qui conditionne la génération sur des propriétés spécifiques (propritétés physico-chimiques, activité biologique)
    \end{itemize}
    \end{block}
\end{frame}

\begin{frame}
    \frametitle{Résultats du générateur}
    \begin{block}{VAE et CVAE}
        \begin{itemize}
        \item Génération de peptides aléatoires à l'aide du VAE
        \end{itemize}
        \centering
        01: IWIGRK

        \begin{itemize}
        \item Génération de peptides à partir de trois propriétés différentes (longueur, cathionicity, hydrophobicity) à l'aide du CVAE
        \end{itemize}

        \centering
        01: KRACLIRRLR

        \begin{center}
        \begin{tabular}{|c|c|}
        \hline
        Lettre & Composition chimique \\
        \hline
        K & Lysine (K) \\
        R & Arginine (R) \\
        A & Alanine (A) \\
        C & Cysteine (C) \\
        L & Leucine (L) \\
        I & Isoleucine (I) \\
        \hline
        \end{tabular}
        \end{center}
    \end{block}
\end{frame}
\begin{frame}
    \begin{block}{Évaluation du générateur}
        \begin{itemize}
        \item Évaluation de la qualité des peptides générés en comparant leurs propriétés à celles des peptides connus
        \end{itemize}
        \centering
        Lorsqu'on donne la sequence KLFKRWKHLFR de taille 11 et ses propriétés physico-chimiques au CVAE, il génère la sequence RKLKLFFRLLR avec des propriétés similaires.

        \begin{itemize}
        \item Utilisation de métriques telles que la diversité, la validité et la nouveauté des séquences générées
        \end{itemize}
    \end{block}
\end{frame}


